\chapter{Introduction to Verilog}

\section{Why Hardware Description Languages Exist}

As digital systems grew beyond small collections of gates, schematic-only design became unsuitable
for serious integrated-circuit development. A modern chip must be described, simulated, revised, and
verified long before it reaches physical implementation, and this requires a representation that EDA
tools can process consistently. \textbf{Hardware Description Languages (HDLs)} provide this
representation. They allow engineers to describe logic, registers, interconnections, timing intent,
and hierarchy in text, while still representing hardware rather than software instructions executed
by a processor.\\

\noindent \textbf{Verilog} is one of the primary HDLs used for digital design and verification. It can
describe simple combinational circuits, sequential elements, counters, multiplexers, and complete
RTL blocks that later pass through simulation, synthesis, and implementation. Its main value in this
flow is that the same design can be organized hierarchically, tested before fabrication, and reused
across larger systems. This chapter introduces the basic Verilog program structure, modules, ports,
data types, constants, and operators needed for the RTL design and verification work used in later
chapters.\\


\section{Verilog Program Structure}

A Verilog source file is organized around one or more modules, and each module acts as a self-contained
hardware block with its own ports, declarations, and internal logic. This structure is not a stylistic
convention; it is what allows EDA tools to parse the design, elaborate the hierarchy, and determine
which parts of the description are meant to behave as combinational logic, sequential logic, or
hierarchical instances of other blocks. A complete Verilog design therefore begins with the module
header, continues with the declaration of ports and internal objects, and ends with one or more
behavioral or structural descriptions that define how the block operates.\\

\noindent The general pattern is consistent across small laboratory examples and larger RTL blocks.
The module header names the design unit and lists its external interface. The body then defines
whether each port is an input, output, or inout, assigns widths where required, declares nets and
variables, and introduces continuous assignments or procedural blocks. A module may also contain
parameters, submodule instances, tasks, functions, and generate constructs when the design grows in
complexity. The following skeleton shows the layout used throughout this manual:

\begin{codeblock}{verilog}
module module_name (port_list);
    // port directions and widths
    // internal declarations
    // continuous assignments
    // procedural blocks
endmodule
\end{codeblock}

\noindent In practice, this structure allows the same description to be read by both simulator and
synthesis tools. Simulation uses the procedural and concurrent statements to model behavior over
time, while synthesis extracts the logic implied by the supported constructs. For that reason, the
organization of the code matters as much as the syntax itself. A clean module layout makes a design
easier to review, easier to debug, and easier to hand over to another engineer or laboratory
administrator later in the flow.

\section{Modules}

The module is the basic unit of composition in Verilog. Every hardware description is built from one
or more modules, whether the design represents a simple gate-level block or a complete subsystem with
multiple submodules. A module may be reused many times within a larger design, which is why it is
useful to think of it as a hardware component rather than as a software procedure. When one module
instantiates another, the relationship is fixed at elaboration time; the resulting hierarchy mirrors
the structure of the hardware that will be simulated or synthesized.\\

\noindent The module body can include instances of other modules, primitive gates, continuous
assignments, tasks, functions, and procedural blocks. A small module may contain only a few lines of
declaration and a single assignment, while a larger module may include several nested blocks and many
internal signals. The important point is that the interface remains explicit. For example, a full
adder can be described with only a few statements:

\begin{codeblock}{verilog}
module full_adder(A, B, Cin, S, Cout);
    input A, B, Cin;
    output S, Cout;

    assign {Cout, S} = A + B + Cin;
endmodule
\end{codeblock}

\noindent The same block may later be reused as part of a larger arithmetic unit without changing its
interface. This reuse is one of the practical reasons Verilog became central to RTL design. It allows
an engineer to define a small verified building block once, then connect that block into a wider
system by instantiating it repeatedly. The module boundary therefore serves both as a design boundary
and as a documentation boundary.

\section{Ports}

Ports define the external interface of a module. They are the connection points through which the
module exchanges signals with the rest of the design, and they must be declared explicitly as input,
output, or inout. A port list alone is not enough to define the signal type; the direction and, when
necessary, the data type and bit range must also be stated inside the module body. In many cases a
port is treated as a wire by default, but wider buses and registered outputs require additional
declaration.\\

\noindent The port list becomes especially important when a module is parameterized. Widths can be
described using parameters so that the same source file may be reused for different data widths
without rewriting the logic. This is common in laboratory designs, where a small example may later be
scaled to a wider bus for experimentation. The syntax below shows the basic pattern:

\begin{codeblock}{verilog}
module module_name
    #(parameter WORD = 32)
    (addr, sum, a, b, sel, data_bus);

    input  [WORD-1:0] addr;
    output reg signed [32:1] sum;
    input  a, b, sel;
    input  [0:15] data_bus;
endmodule
\end{codeblock}

\noindent The declaration order matters when the design is read by tools and by human reviewers.
Signals in the port list should be easy to map back to the declarations, and the direction of each
signal should be obvious before the body of the module is examined. A clear port definition is one of
the simplest ways to keep a Verilog project maintainable as it grows.

\section{Nets and Registers}

Verilog separates physical connections from storage elements. Nets, such as \inlinecode{wire}, are
used to model connections between components, while registers such as \inlinecode{reg} are used when
the value must be held or assigned inside procedural code. This distinction is important because it
reflects the hardware interpretation of the source. A net is driven by some source and used as a
connection in the design, whereas a register is typically updated by an \inlinecode{always} or
\inlinecode{initial} block during simulation.\\

\noindent Verilog also models four-valued logic, which is necessary when a design is not fully driven
or when a signal cannot yet be resolved during simulation. The values are \inlinecode{0},
\inlinecode{1}, \inlinecode{x}, and \inlinecode{z}. The \inlinecode{z} state represents a high-
impedance or undriven condition, while \inlinecode{x} represents an unknown or unresolved value.
These states are useful when checking reset behavior, tri-state buses, or incomplete initialization.
The following examples show the basic declaration style used in practice:

\begin{codeblock}{verilog}
wire A, X, Y, Z, Cin, Cout;
wire [3:0] B, S;
wire [15:0] C;

reg q_out;
reg [7:0] counter;
\end{codeblock}

\noindent Nets are not assigned inside procedural blocks, and registers are not used as direct
continuous-connection targets in the same way as a wire. That separation keeps the model consistent
with the intended hardware behavior and helps the simulator distinguish between continuous hardware
connections and values that are updated under procedural control.

\section{Numbers and Constants}

Verilog numbers are written using a compact notation that encodes the size, radix, and value in a
single literal. The general form is \inlinecode{<size>'<radix><value>}, where the radix indicates
binary, octal, decimal, or hexadecimal representation. This syntax is used throughout RTL coding
because it makes the intended width explicit and avoids ambiguity about sign extension or default
integer size. The standard radix characters are \inlinecode{b}, \inlinecode{o}, \inlinecode{d}, and
\inlinecode{h}.\\

\noindent Sized constants are preferable whenever the width matters, especially in module interfaces,
state machines, and testbench stimulus. Verilog also allows unknown and high-impedance digits inside
constants, which is useful when modeling incomplete bus values or uninitialized conditions during
simulation. The examples below match the style used in the notes:

\begin{codeblock}{verilog}
4'b1010
8'h3F
16'd255
4'b01xx
1'bz
\end{codeblock}

\noindent Constants are not simply a syntactic detail. In a hardware description, they determine how
the simulator interprets a value and how a synthesis tool treats a mask, width, or reset state. A
value written without an explicit size may behave differently from the same value written with a
fixed bit width, so laboratory code should state widths clearly whenever the design depends on them.

\section{Operators}

Verilog operators follow the same broad categories used in digital logic, but they are applied
directly to hardware signals rather than to software variables. Bitwise operators act on individual
bits, logical operators act on the truth value of an expression, shift operators move data left or
right, and the conditional operator provides a compact form of multiplexing. Understanding this
division is essential because a mistake in operator choice can change the synthesized hardware even
when the simulation appears to run.\\

\noindent Bitwise operators operate on each bit of their operands. Unary reduction operators, by
contrast, collapse a vector into a single result by combining all bits together. Logical operators
return a true or false result and are typically used in conditions. The shift operators
\inlinecode{<<} and \inlinecode{>>} move the bit pattern by a fixed number of positions, while the
conditional operator \inlinecode{?:} is commonly used to express a multiplexer in a compact form.
The basic usage is illustrated below:

\begin{codeblock}{verilog}
assign A = X | (Y & ~Z);
assign B = &data_bus;
assign C = |mask_bus;
assign D = sel ? in1 : in0;

a = 4'b1010;
d = a >> 2;
c = a << 1;
\end{codeblock}

\noindent In laboratory RTL, these operators are used constantly because they map directly onto the
hardware view of the design. A conditional expression often describes the same logic that would be
drawn as a multiplexer in a schematic, while bitwise and reduction operators are useful for masks,
parity checks, and grouped control logic. The operator chosen in the source should therefore match
the intended hardware structure, not merely the result visible in a quick simulation.

\section{Laboratory Exercise -- Half Adder}

\subsection{Objective}
The purpose of this exercise is to design, simulate, and verify a Half Adder using Verilog HDL. The
exercise introduces a complete but minimal digital design flow, beginning with the RTL description
and continuing through testbench development, compilation, simulation, and waveform inspection in
the Cadence environment.\\

\noindent A Half Adder is a suitable starting point because it is small enough to understand in a
single session, yet complete enough to demonstrate the relationship between Boolean equations,
Verilog code, stimulus generation, and simulation results. The same workflow used here will be used
again in later chapters for larger combinational and sequential designs.

\subsection{Theory}
A Half Adder is a combinational circuit that adds two one-bit inputs, \inlinecode{A} and
\inlinecode{B}, and produces two one-bit outputs: the \textit{Sum} and the \textit{Carry}. The Sum
bit represents the least significant result of the addition, while the Carry bit indicates whether
both inputs were high at the same time. Since the circuit has no storage element, its outputs depend
only on the current input combination.\\


\begin{wrapfigure}{r}{0.3\textwidth}
	\centering
	\includegraphics[width=0.35\textwidth]{resources/half_adder.jpg}
\end{wrapfigure}
\noindent In Boolean form, the Half Adder is expressed as:
\[
\text{Sum} = A \oplus B
\]
\[
\text{Carry} = A \cdot B
\]

\noindent The Sum output is obtained using the exclusive-OR operation because the result is high
only when exactly one input is high. The Carry output is obtained using the AND operation because a
carry occurs only when both inputs are asserted simultaneously.


\subsection{Truth Table}
The truth table below lists all possible input combinations for a Half Adder and the corresponding
Sum and Carry outputs. This table serves as the reference against which the Verilog model and the
testbench results are checked.

\begin{table}[H]
\centering
\renewcommand{\arraystretch}{1.2}
\begin{tabular}{|c|c|c|c|}
\hline
\textbf{A} & \textbf{B} & \textbf{Sum} & \textbf{Carry} \\
\hline
0 & 0 & 0 & 0 \\
\hline
0 & 1 & 1 & 0 \\
\hline
1 & 0 & 1 & 0 \\
\hline
1 & 1 & 0 & 1 \\
\hline
\end{tabular}
\caption{Truth table of the Half Adder.}
\label{tab:half_adder_truth}
\end{table}

\subsection{Verilog Implementation}
The Verilog description of a Half Adder is intentionally short, but it should still follow the same
discipline used for larger RTL blocks. The module name identifies the design unit, the port list
exposes the external interface, and the continuous assignments map directly to the Boolean
expressions derived from the truth table. This makes the source easy to read and easy to verify
against the expected hardware behavior.\\

\begin{codeblock}{verilog}
module half_adder(A, B, Sum, Carry);
    input A, B;
    output Sum, Carry;

    assign Sum = A ^ B;
    assign Carry = A & B;
endmodule
\end{codeblock}

\noindent The module declaration names the design and defines the port order. The input declarations
identify the two one-bit operands. The output declarations identify the result signals that will be
driven by the logic inside the module. The \inlinecode{assign} statement for \inlinecode{Sum}
implements the exclusive-OR function, while the \inlinecode{assign} statement for \inlinecode{Carry}
implements the AND function. The \inlinecode{endmodule} keyword closes the design unit and allows it
to be instantiated inside a testbench or a larger circuit.

\subsection{Testbench}
A testbench is required because the design itself does not generate stimuli. The Half Adder must be
exercised with all valid input combinations so that the simulation can confirm both outputs against
the truth table. In a laboratory setting, the testbench also establishes a repeatable procedure for
running the design inside the simulator, which is essential when the same flow must later be repeated
for larger circuits.\\

\begin{codeblock}{verilog}
`timescale 1ns/1ps

module half_adder_tb;
    reg A, B;
    wire Sum, Carry;

    half_adder dut (
        .A(A),
        .B(B),
        .Sum(Sum),
        .Carry(Carry)
    );

    initial begin
        $monitor($time, " A=%b B=%b Sum=%b Carry=%b", A, B, Sum, Carry);

        A = 0; B = 0; #10;
        A = 0; B = 1; #10;
        A = 1; B = 0; #10;
        A = 1; B = 1; #10;

        $finish;
    end
endmodule
\end{codeblock}

\noindent The stimulus sequence applies each of the four possible input combinations in turn, with a
short delay between each state so that the waveform viewer can capture the transitions clearly. The
\inlinecode{$monitor} statement records the active values during simulation, making it easy to compare
the observed output against the truth table. Once the final input combination has been applied, the
simulation terminates with \inlinecode{$finish}.

\subsection{Simulation Procedure}
The simulation workflow should be treated as a standard laboratory procedure rather than as an ad hoc
sequence of tool clicks. A clean project directory, a correct set of source files, and a known
simulation entry point make it much easier to reproduce results later. The exact window layout may
vary slightly between Cadence releases, but the workflow remains the same.\\

\begin{enumerate}
    \item Create a new Cadence project or working directory for the exercise.
    \item Add the design file containing the Half Adder RTL description.
    \item Add the testbench file and ensure it references the design module correctly.
    \item Compile the design and confirm that no syntax errors are reported.
    \item Run the simulation and allow the testbench to apply all input combinations.
    \item Open the waveform viewer and inspect the Sum and Carry signals.
    \item Save the simulation session if the laboratory workflow requires a record of the results.
\end{enumerate}

% TODO: Insert compilation screenshot
% TODO: Insert waveform viewer screenshot

\noindent In practice, the compilation stage confirms that both files are readable by the simulator
and that the module names, port names, and signal declarations agree. The simulation stage then
checks that the testbench drives the design in the intended order. The waveform viewer is used at the
end because it provides the clearest visual confirmation that the output transitions match the truth
table.

\subsection{Expected Results}
The expected waveform is straightforward. When both inputs are low, both outputs remain low. When
exactly one input is high, the Sum output becomes high while the Carry output remains low. When both
inputs are high, the Sum output returns low and the Carry output becomes high. This is the defining
behavior of a Half Adder and should be visible directly in the waveform viewer.\\

\begin{figure}[H]
    \centering
    \fbox{\parbox{0.82\linewidth}{\centering TODO: Insert Half Adder waveform screenshot}}
    \caption{Expected Half Adder waveform showing Sum and Carry behavior for all input combinations.}
    \label{fig:half_adder_waveform}
\end{figure}

\noindent The waveform should therefore show \inlinecode{Sum} following the XOR of \inlinecode{A} and
\inlinecode{B}, while \inlinecode{Carry} follows the AND of the same inputs. If the observed outputs
do not match the truth table, the design file and testbench should be checked first, followed by the
simulation setup and signal connections.

\subsection{Conclusion}
This exercise demonstrates the complete Verilog verification flow on a simple combinational circuit.
The Half Adder is small, but it introduces the essential sequence used throughout digital design:
derive the logic, write the RTL, apply test stimulus, simulate the model, and verify the waveform
against the expected behavior. That procedure will be reused in later chapters when the designs
become larger and the simulation setup becomes more involved.\\

\noindent By completing the exercise, the reader gains a practical starting point for working with
Cadence tools in a reproducible laboratory environment. The same method of checking the Boolean
equations, validating the truth table, and confirming the waveform response will remain useful when
the designs move from simple combinational logic to sequential circuits and larger RTL blocks.
